В соответствии с задачами курсовой работы был проведён обзор классических типов аллокаторов: линейного, пуллового и стекового. Был также выполнен анализ их быстродействия по сравнению со стандартным аллокатором. Для этого потребовалось реализовать каждый из них на языке C++. При реализации был выбран путь максимального упрощения, избегания сложных структур данных (односвязный список и стек, потребовавшиеся в данных аллокаторах, были реализованы самостоятельно) с целью достижения максимальной эффективности.

Помимо этого было проанализировано хранение в памяти не только простых объектов, но и структур. На основе идеи колоночных баз данных был разработан аллокатор, который хранит данные структур не по строкам, а по столбцам, и было проведено его сравнение по скорости выделения и освобождения памяти с обычным пулловым аллокатором.

В результате исследования были сделаны следующие выводы:

\begin{itemize}
\item C++ "--- язык с ручным управлением памятью. С одной стороны, это обеспечивает гибкость и универсальность при работе с памятью. С другой стороны, это является причиной таких проблем, как фрагментация и утечки памяти, внезапные зависания и снижение производительности. Многих вышеперечисленных проблем можно избежать, если разработать собственный аллокатор.
\item Линейный, стековый и пулловый аллокатор являются достаточно простыми в реализации, в то же время их выигрыш в производительности по сравнению со стандартным аллокатором огромен: было достигнуто ускорение в $3,25$ (для пуллового) -- $19,05$ (для линейного и стекового) раз. Тем не менее, стоит упомянуть, что, как правило, операторы \texttt{new} и \texttt{delete} не используются при создании простых объектов: вместо этого можно было бы создать массив или вектор чисел, которые, вероятно, работали бы не медленнее, чем пользовательские аллокаторы.
\item Необходимо избегать бездумных вызовов \texttt{new} и \texttt{delete}, и заранее выделять память под большие массивы объектов, чтобы ускорить работу программы. Это особенно важно в тех ситуациях, где многое зависит от производительности: в играх, в высоконагруженных приложениях, в олимпиадном программировании.
\item Реализованные в работе аллокаторы не являются универсальными: линейный и стековый аллокатор недостаточно функциональны в вопросах, касающихся освобождения памяти, а пулловый аллокатор приспособлен только под выделение блоков памяти фиксированного размера. Если необходимо разработать универсальный собственный аллокатор, можно изучить такие стратегии, как <<First fit>> или <<Buddy Allocator>> \cite{memory-management}.
\end{itemize}

Далеко не у каждого разработчика хотя бы раз в жизни возникает необходимость реализовать собственный аллокатор, и для большинства нужд вполне подходящим оказывается стандартный аллокатор языка C++. Тем не менее, изучение различных алгоритмов аллокации позволяет лучше понять то, как происходит ручная работа с памятью, выявить причины замедлений в программе и, в случае обнаружения проблем, связанных с управлением памятью, заменить стандартный алгоритм аллокации некоторым пользовательским вариантом.
